% Source: https://artofproblemsolving.com/community/c3h2949415_2021_emcc_guts_round_p13_p24_exeter_math_club_competition
% Stable identity: aops-topic-2949415-q18
Each region consists of its fixed circular segment and a triangle with the corresponding square side as base. Opposite triangles have equal bases and heights summing to the distance between the opposite sides, so the two opposite regions sum to half the disk. The other fractions are consequently $\dfrac{1}{2}-\dfrac{1}{5}=\dfrac{3}{10}$ and $\dfrac{1}{2}-\dfrac{1}{3}=\dfrac{1}{6}$. In descending order these are $\dfrac{1}{3},\dfrac{3}{10},\dfrac{1}{5},\dfrac{1}{6}$, so the answer is $\boxed{\dfrac{3}{10}}$.
\begin{asy}
usepackage("amsmath");
size(7cm);
pair A=(1,0),B=(0,1),C=(-1,0),D=(0,-1),P=(2pi/15,-pi/30);draw(circle((0,0),1));draw(A--P--B^^C--P--D);label("$A$",A,E);label("$B$",B,N);label("$C$",C,W);label("$D$",D,S);label("$P$",P,NE);
\end{asy}
